\documentclass[11pt,a4paper]{article}
\usepackage[margin=2.5cm]{geometry}
\usepackage{amsmath,amssymb}

\newcommand{\R}{\mathbb{R}}
\newcommand{\E}{\mathbb{E}}
\newcommand{\norm}[1]{\left\| #1 \right\|}
\newcommand{\ip}[2]{\left\langle #1, #2 \right\rangle}

\title{Homework 19.10.26 -- Exercises 20, 22--26}
\author{Gor Melkumyan}
\date{}

\begin{document}
\maketitle

\section*{Exercise 20}

Given that $f$ has $L$-Lipschitz gradient, i.e.,
\[
  \norm{\nabla f(x) - \nabla f(y)} \le L\norm{x-y} \qquad \forall x, y \in \R^d .
\]
Show that
\[
  -\frac{L}{2}\norm{x-y}^2 \le D_f(x,y) \le \frac{L}{2}\norm{x-y}^2 \qquad \forall x, y \in \R^d .
\]

\subsection*{Proof}
It is enough to show that $|D_f(x,y)| \le \frac{L}{2}\norm{x-y}^2$.
Let $z := y + t(x-y)$, so that $x - y = \frac{z-y}{t}$. Then
\begin{align*}
  |D_f(x,y)|
  &= \bigl|f(x) - f(y) - \ip{\nabla f(y)}{x-y}\bigr|
  \\
  &= \left|\int_0^1 \ip{\nabla f\bigl(y + t(x-y)\bigr)}{x-y}\,dt - \ip{\nabla f(y)}{x-y}\right| \\
  &= \left|\int_0^1 \ip{\nabla f(z) - \nabla f(y)}{z-y}\cdot\frac{1}{t}\,dt\right|
  \le \int_0^1 \frac{1}{t}\,\bigl|\ip{\nabla f(z) - \nabla f(y)}{z-y}\bigr|\,dt .
\end{align*}
By Cauchy--Schwarz and then by the $L$-Lipschitz gradient assumption,
\begin{align*}
  |D_f(x,y)|
  &\le \int_0^1 \frac{1}{t}\norm{z-y}\norm{\nabla f(z) - \nabla f(y)}\,dt
  \le \int_0^1 \frac{L}{t}\norm{z-y}^2\,dt
  = \int_0^1 \frac{L}{t}\ip{z-y}{z-y}\,dt \\
  &= \int_0^1 L\ip{x-y}{x-y}\,t\,dt
  = \frac{L}{2}\norm{x-y}^2,
\end{align*}
where we used $z - y = t(x-y)$.

\section*{Exercise 22}

Let $T_\gamma(x) := x - \gamma\nabla f(x)$. Show that if $f$ has $L$-Lipschitz gradient, then
for all $x \in \R^d$ and $\gamma \ge 0$
\[
  \gamma\left(1 - \frac{\gamma L}{2}\right)\norm{\nabla f(x)}^2
  \le f(x) - f\bigl(T_\gamma(x)\bigr)
  \le \gamma\left(1 + \frac{3\gamma L}{2}\right)\norm{\nabla f(x)}^2 .
\]

\subsection*{Proof}
Let $x_t := x - t\gamma\nabla f(x)$. From the Fundamental Theorem of Calculus we get
\begin{align*}
  f(x) - f\bigl(T_\gamma(x)\bigr)
  &= \int_0^1 \ip{\nabla f\Bigl(x - t\bigl(x - (x - \gamma\nabla f(x))\bigr)\Bigr)}{x - \bigl(x - \gamma\nabla f(x)\bigr)}\,dt \\
  &= \int_0^1 \ip{\nabla f\bigl(x - t\gamma\nabla f(x)\bigr)}{\gamma\nabla f(x)}\,dt
  = \gamma\int_0^1 \ip{\nabla f(x)}{\nabla f(x_t)}\,dt .
\end{align*}
Also we have
\[
  \ip{\nabla f(x)}{\nabla f(x_t)} = \norm{\nabla f(x)}^2 + \ip{\nabla f(x)}{\nabla f(x_t) - \nabla f(x)} .
\]
As a result,
\[
  f(x) - f\bigl(T_\gamma(x)\bigr)
  = \gamma\norm{\nabla f(x)}^2 + \gamma\int_0^1 \ip{\nabla f(x)}{\nabla f(x_t) - \nabla f(x)}\,dt .
\]
So if we prove that
\[
  -\frac{\gamma L}{2}\norm{\nabla f(x)}^2
  \le \int_0^1 \ip{\nabla f(x)}{\nabla f(x_t) - \nabla f(x)}\,dt
  \le \frac{3\gamma L}{2}\norm{\nabla f(x)}^2,
\]
then the statement follows.
On the other hand, by Cauchy--Schwarz and the $L$-Lipschitz gradient,
\begin{align*}
  \bigl|\ip{\nabla f(x)}{\nabla f(x_t) - \nabla f(x)}\bigr|
  &\le \norm{\nabla f(x)}\cdot\norm{\nabla f(x_t) - \nabla f(x)}
  \le L\norm{\nabla f(x)}\cdot\norm{x_t - x} \\
  &= L\norm{\nabla f(x)}\cdot\norm{t\gamma\nabla f(x)}
  = L\gamma t\norm{\nabla f(x)}^2 .
\end{align*}
As a result,
\[
  \left|\int_0^1 \ip{\nabla f(x)}{\nabla f(x_t) - \nabla f(x)}\,dt\right|
  \le L\gamma\norm{\nabla f(x)}^2\int_0^1 t\,dt
  = \frac{L\gamma}{2}\norm{\nabla f(x)}^2,
\]
which proves the statement, because $\frac{L\gamma}{2} \le \frac{3L\gamma}{2}$.

\section*{Exercise 23}

Find an upper bound for $\norm{x^{k+1} - x^*}^2$ using Proposition 20, i.e.: if
$f\colon \R^d \to \R$, $f \in C^1(\R^d)$, is also $\mu$-convex and $L$-smooth, then
\[
  \mu\norm{x-y}^2 + \frac{1}{L}\norm{\nabla f(x) - \nabla f(y)}^2
  \le \left(1 + \frac{\mu}{L}\right)\ip{\nabla f(x) - \nabla f(y)}{x-y}
  \qquad \forall x, y \in \R^d .
\]

\subsection*{Proof}
\begin{align*}
  \norm{x^{k+1} - x^*}^2
  &= \norm{(x^k - x^*) - \gamma\bigl(\nabla f(x^k) - \nabla f(x^*)\bigr)}^2 \\
  &= \norm{x^k - x^*}^2 + \gamma^2\norm{\nabla f(x^k) - \nabla f(x^*)}^2
     - 2\gamma\ip{\nabla f(x^k) - \nabla f(x^*)}{x^k - x^*} .
\end{align*}
On the other hand, multiplying Proposition 20 by $\frac{L}{L+\mu}$,
\[
  \ip{\nabla f(x^k) - \nabla f(x^*)}{x^k - x^*}
  \ge \frac{L\mu}{L+\mu}\norm{x^k - x^*}^2 + \frac{1}{L+\mu}\norm{\nabla f(x^k) - \nabla f(x^*)}^2 .
\]
Multiplying by $-2\gamma$:
\[
  -2\gamma\ip{\nabla f(x^k) - \nabla f(x^*)}{x^k - x^*}
  \le -\frac{2L\mu\gamma}{L+\mu}\norm{x^k - x^*}^2 - \frac{2\gamma}{L+\mu}\norm{\nabla f(x^k) - \nabla f(x^*)}^2 .
\]
As a result,
\[
  \norm{x^{k+1} - x^*}^2
  \le \left(1 - \frac{2L\mu\gamma}{L+\mu}\right)\norm{x^k - x^*}^2
  - \gamma\left(\frac{2}{L+\mu} - \gamma\right)\norm{\nabla f(x^k) - \nabla f(x^*)}^2 .
\]
In particular, for $0 \le \gamma \le \frac{2}{L+\mu}$ the last term is non-positive, and we get
\[
  \norm{x^{k+1} - x^*}^2 \le \left(1 - \frac{2L\mu\gamma}{L+\mu}\right)\norm{x^k - x^*}^2 .
\]
Note that
\[
  1 - \frac{2L\mu\gamma}{L+\mu} \le 1 - \mu\gamma,
  \qquad \text{because} \qquad \frac{2L}{L+\mu} \ge 1 .
\]

\section*{Exercise 24}

For all $a, b \in \R^d$ and $t > 0$ the following inequalities hold:
\begin{align*}
  &1)\quad \ip{a}{b} \le \frac{\norm{a}^2}{2t} + \frac{t\norm{b}^2}{2}, \\
  &2)\quad \norm{a+b}^2 \le 2\norm{a}^2 + 2\norm{b}^2, \\
  &3)\quad \frac{1}{2}\norm{a}^2 - \norm{b}^2 \le \norm{a+b}^2 .
\end{align*}

\subsection*{Proof of 1)}
Let $a' = \frac{a}{\sqrt{t}}$, $b' = b\sqrt{t}$, i.e., $a = a'\sqrt{t}$, $b = \frac{b'}{\sqrt{t}}$. Then
\[
  \ip{a}{b} = \ip{a'}{b'}, \qquad
  \frac{\norm{a}^2}{2t} = \frac{\norm{a'}^2}{2}, \qquad
  \frac{t\norm{b}^2}{2} = \frac{\norm{b'}^2}{2}.
\]
As a result, we only have to prove
\[
  2\ip{a'}{b'} \le \norm{a'}^2 + \norm{b'}^2,
\]
which is equivalent to $\norm{a'-b'}^2 \ge 0$, so it is true.

\subsection*{Proof of 2)}
Using $\norm{a-b}^2 \ge 0$ we get
\[
  \norm{a}^2 + \norm{b}^2 - 2\ip{a}{b} \ge 0 .
\]
Adding $\norm{a}^2 + \norm{b}^2$ to both sides:
\[
  2\norm{a}^2 + 2\norm{b}^2 - 2\ip{a}{b} \ge \norm{a}^2 + \norm{b}^2 .
\]
Adding $2\ip{a}{b}$ and regrouping:
\[
  2\norm{a}^2 + 2\norm{b}^2 \ge \norm{a}^2 + 2\ip{a}{b} + \norm{b}^2 = \norm{a+b}^2 .
\]

\subsection*{Proof of 3)}
3) is the same as
\[
  \norm{a}^2 \le 2\norm{a+b}^2 + 2\norm{b}^2 .
\]
Setting $a := a' + b'$, $b := -b'$ (so that $a + b = a'$) we get
\[
  \norm{a'+b'}^2 \le 2\norm{a'}^2 + 2\norm{b'}^2,
\]
which is the same as 2).

\section*{Exercise 25}

For a random vector $X \in \R^d$ (with finite second moment) and any $c \in \R^d$,
\[
  \E\bigl[\norm{X - \E[X]}^2\bigr] = \E\bigl[\norm{X - c}^2\bigr] - \norm{\E[X] - c}^2 .
\]

\subsection*{Proof}
\begin{align*}
  \E\bigl[\norm{X - \E[X]}^2\bigr]
  &= \E\bigl[\norm{(X - c) - (\E[X] - c)}^2\bigr] \\
  &= \E\bigl[\norm{X - c}^2\bigr] + \E\bigl[\norm{\E[X] - c}^2\bigr]
     - 2\,\E\bigl[\ip{X-c}{\E[X]-c}\bigr].
\end{align*}
On the other hand, $\norm{\E[X] - c}^2$ is a number, so its expectation is the same number:
\[
  \E\bigl[\norm{\E[X] - c}^2\bigr] = \norm{\E[X] - c}^2 .
\]
So we get
\[
  \E\bigl[\norm{X - \E[X]}^2\bigr]
  = \E\bigl[\norm{X - c}^2\bigr] - \norm{\E[X] - c}^2
    - 2\Bigl(\E\bigl[\ip{X-c}{\E[X]-c}\bigr] - \norm{\E[X] - c}^2\Bigr).
\]
Also
\begin{align*}
  \E\bigl[\ip{X-c}{\E[X]-c}\bigr]
  &= \E\bigl[\ip{X}{\E[X]}\bigr] - \ip{c}{\E[X]} - \ip{c}{\E[X]} + \norm{c}^2 \\
  &= \norm{\E[X]}^2 - 2\ip{c}{\E[X]} + \norm{c}^2 = \norm{\E[X] - c}^2,
\end{align*}
which means
\[
  -2\Bigl(\E\bigl[\ip{X-c}{\E[X]-c}\bigr] - \norm{\E[X] - c}^2\Bigr) = 0,
\]
which proves the identity.


\section*{Exercise 26}

Prove that for a random vector $X \in \R^d$ and any $t > 0$ we have
\[
  \mathrm{Prob}\bigl(\norm{X - \E[X]} \ge t\bigr) \le \frac{\mathrm{Var}[X]}{t^2} .
\]
Below $\mu := \E[X]$ and $\sigma^2 := \mathrm{Var}[X] = \E\bigl[\norm{X - \mu}^2\bigr]$.

\subsection*{Proof 1 (from scratch)}
Replacing $t$ by $t\sigma$ (for $\sigma > 0$),
\[
  \mathrm{Prob}\bigl(\norm{X - \mu} \ge t\bigr) \le \frac{\sigma^2}{t^2}
  \iff
  \mathrm{Prob}\bigl(\norm{X - \mu} \ge t\sigma\bigr) \le \frac{1}{t^2} .
\]
By the law of total expectation,
\begin{align*}
  \sigma^2 = \E\bigl[\norm{X - \mu}^2\bigr]
  &= \E\bigl[\norm{X - \mu}^2 \,\big|\, \norm{X - \mu} \ge t\sigma\bigr]\,
     \mathrm{Prob}\bigl(\norm{X - \mu} \ge t\sigma\bigr) \\
  &\quad + \E\bigl[\norm{X - \mu}^2 \,\big|\, \norm{X - \mu} < t\sigma\bigr]\,
     \mathrm{Prob}\bigl(\norm{X - \mu} < t\sigma\bigr) \\
  &\ge (t\sigma)^2\,\mathrm{Prob}\bigl(\norm{X - \mu} \ge t\sigma\bigr)
     + 0\cdot\mathrm{Prob}\bigl(\norm{X - \mu} < t\sigma\bigr) \\
  &= (t\sigma)^2\,\mathrm{Prob}\bigl(\norm{X - \mu} \ge t\sigma\bigr).
\end{align*}
As a result, $(t\sigma)^2\,\mathrm{Prob}\bigl(\norm{X - \mu} \ge t\sigma\bigr) \le \sigma^2$, which is the same as
\[
  \mathrm{Prob}\bigl(\norm{X - \mu} \ge t\sigma\bigr) \le \frac{1}{t^2} .
\]

\subsection*{Proof 2 (using Markov's inequality)}
Markov's inequality: for a non-negative random variable $Z$ and $s > 0$,
\[
  \mathrm{Prob}(Z \ge s) \le \frac{\E[Z]}{s} .
\]
Let $Z := \norm{X - \mu}^2$, where $\mu = \E[X]$. Then
\[
  \mathrm{Prob}\bigl(\norm{X - \mu}^2 \ge s\bigr) \le \frac{\E\bigl[\norm{X - \mu}^2\bigr]}{s},
\]
which is the same as
\[
  \mathrm{Prob}\bigl(\norm{X - \mu} \ge \sqrt{s}\bigr) \le \frac{\mathrm{Var}[X]}{s} .
\]
Substituting $s = t^2$ we get
\[
  \mathrm{Prob}\bigl(\norm{X - \mu} \ge t\bigr) \le \frac{\mathrm{Var}[X]}{t^2},
\]
which proves the inequality.

\end{document}
